% Appendix: GED on the base witness.

\section{Generalised Euclidean distance on the base witness}\label{app:ged-witness}

This appendix computes the generalised Euclidean distance
(Definition~\ref{def:ged}) on the pair $(A, B) = (C_3 \sqcup C_3, C_6)$ of
Section~\ref{sec:warmup} under an opinion vector aligned with the natural
bipartition of each graph. The computation establishes part (d) of
Proposition~\ref{prop:noncomputable-measures}.

Vertices are indexed $1$ through $6$. Take $S_+ = \{1, 2, 3\}$ and
$S_- = \{4, 5, 6\}$. On $A$, $S_+$ is the first triangle and $S_-$
the second. On $B$ we use the bipartition into $\{1, 2, 3\}$ and
$\{4, 5, 6\}$ (three consecutive vertices in each set); this is not
the alternating $2$-colouring of $C_6$ but does align with two arcs
of the cycle.

Take the opinion vector $o = (+1, +1, +1, -1, -1, -1)^\top$. Its
positive and negative parts are $o^+ = (1, 1, 1, 0, 0, 0)^\top$ and
$o^- = (0, 0, 0, 1, 1, 1)^\top$. Their difference is $o^+ - o^- = o$.

The expansion given by \citet{HohmannGED2023} for the squared
distance, treating $o^+$ and $o^-$ as probability distributions after
normalisation, is
\[
\delta_{G,o}^2 = \mathbb{E}[\omega_{X^+, X^-}] - \tfrac{1}{2}\mathbb{E}[\omega_{X^+, X^+}] - \tfrac{1}{2}\mathbb{E}[\omega_{X^-, X^-}],
\]
where $X^+$ is drawn from the normalised $o^+$, $X^-$ is drawn from
the normalised $o^-$ (both are uniform on their respective supports
here), and $\omega_{uv}$ is the effective resistance between $u$ and
$v$ on $\gph$.

\subsection*{Computation on $A$}

$A$ has two connected components, $S_+$ (a triangle) and $S_-$ (a
triangle). Effective resistance across components is not defined on
the graph as a whole. Following the paper's per-component convention,
the pseudoinverse $L^\dagger$ decomposes as the direct sum of
$L_{C_3}^\dagger$ on each component. Concretely,
\[
L_{C_3}^\dagger = \frac{1}{9}(2 I_3 - J_3),
\]
where $I_3$ is the $3\times 3$ identity and $J_3$ is the $3 \times 3$
all-ones matrix. The vector $(1, 1, 1)^\top$ lies in the null space
of $L_{C_3}$ (the space spanned by the all-ones vector), so
$L_{C_3}^\dagger (1, 1, 1)^\top = (0, 0, 0)^\top$.

The quadratic form on the featured graph $A$ is
\[
o^\top L_A^\dagger o = (o|_{S_+})^\top L_{C_3}^\dagger (o|_{S_+}) +
(o|_{S_-})^\top L_{C_3}^\dagger (o|_{S_-}) = 0 + 0 = 0,
\]
where $o|_{S_+} = (1, 1, 1)^\top$ and $o|_{S_-} = (-1, -1, -1)^\top$.
Hence $\mathrm{GED}(A, o) = 0$.

\subsection*{Computation on $B$}

$B = C_6$ is connected. Effective resistance on a cycle $C_n$ between
two vertices at cyclic distance $d$ is $d(n - d) / n$. For $n = 6$:
distance $1$ gives $5/6$, distance $2$ gives $4/3$, distance $3$ gives
$3/2$.

$S_+ = \{1, 2, 3\}$ and $S_- = \{4, 5, 6\}$. Cross-side distances
(cyclic distance from each vertex of $S_+$ to each vertex of $S_-$)
are enumerated in the following table.

\begin{center}
\begin{tabular}{c|ccc}
& $4$ & $5$ & $6$ \\
\hline
$1$ & $3$ & $2$ & $1$ \\
$2$ & $2$ & $3$ & $2$ \\
$3$ & $1$ & $2$ & $3$
\end{tabular}
\end{center}

The nine cross-side effective resistances are $\{3/2, 4/3, 5/6, 4/3,
3/2, 4/3, 5/6, 4/3, 3/2\}$. Their mean is
\[
\mathbb{E}[\omega_{X^+, X^-}] = \frac{1}{9}\left(3 \cdot \tfrac{3}{2}
+ 4 \cdot \tfrac{4}{3} + 2 \cdot \tfrac{5}{6}\right) = \frac{1}{9}\left(\tfrac{9}{2} + \tfrac{16}{3} + \tfrac{5}{3}\right)
= \frac{1}{9} \cdot \frac{27 + 32 + 10}{6} = \frac{69}{54} = \frac{23}{18}.
\]

Within-side distances: for $S_+ = \{1, 2, 3\}$, the three ordered
pairs $(1,1), (2,2), (3,3)$ contribute $0$; the six pairs $(1,2),
(2,1), (2,3), (3,2), (1,3), (3,1)$ contribute distances $1, 1, 1, 1,
2, 2$, giving effective resistances $5/6, 5/6, 5/6, 5/6, 4/3, 4/3$.
The mean over the nine ordered pairs is
\[
\mathbb{E}[\omega_{X^+, X^+}] = \frac{1}{9}\left(4 \cdot \tfrac{5}{6} + 2 \cdot \tfrac{4}{3}\right) = \frac{1}{9}\left(\tfrac{10}{3} + \tfrac{8}{3}\right) = \frac{18}{27} = \frac{2}{3}.
\]
By symmetry, $\mathbb{E}[\omega_{X^-, X^-}] = 2/3$.

Substituting into the paper's expansion:
\[
\mathrm{GED}(B, o)^2 = \frac{23}{18} - \frac{1}{2}\cdot\frac{2}{3} -
\frac{1}{2}\cdot\frac{2}{3} = \frac{23}{18} - \frac{2}{3} = \frac{23}{18} - \frac{12}{18} = \frac{11}{18}.
\]
Hence $\mathrm{GED}(B, o) = \sqrt{11/18} \approx 0.782$.

\subsection*{Consequence}

Under the opinion vector $o = (+1, +1, +1, -1, -1, -1)^\top$ aligned
with the natural bipartition of each graph, $\mathrm{GED}(A, o) = 0$
while $\mathrm{GED}(B, o) = \sqrt{11/18} > 0$. The two graphs are
$1$-WL equivalent (Section~\ref{sec:warmup}). Hence $\mathrm{GED}$ takes
different values on featured graphs that are $1$-WL equivalent, so
$\mathrm{GED}$ is not $1$-WL invariant, as claimed in
Proposition~\ref{prop:noncomputable-measures}(d).
